% Provenance (never printed): see sources/outline_sources.md (Year 2: tholeiitic, calc-alkaline
% and alkaline series; the TAS diagram; settings; isotopic signatures of mantle and crust).
% Classification from BGS RR 99-06 (after Streckeisen 1976, Le Bas et al. 1986, Le Maitre et
% al. 1989); series and data from Vermeesch and Pease (2021) and the GeoplotR package;
% curves and counts from figdata/bachelor-2/series.py; exercise and problem values in the
% Book 4 scratch calculators calc/ch09.py and calc/ch09pb.py.
\chapter{Seri Magmatik dan Tatanan Tektonik}\label{ch:b2:magmatic-series}

Berikan sebuah analisis kimia kepada seorang ahli petrologi tanpa peta, dan ia sering dapat mengatakan apakah
lava itu meletus di punggungan, di pulau samudra, atau di atas lempeng yang menunjam. Lava
Islandia dan lava gunung-gunung api Cascade di bagian barat Amerika Utara sama-sama berkisar
dari basal sampai riolit; bila dianalisis ratusan demi ratusan, keduanya mengikuti dua jalur yang berbeda selama
berevolusi, yang satu makin kaya akan besi, yang lain tidak. Bab ini mula-mula
menamai batuan beku, dari mineralnya bila mineral itu dapat dilihat dan dari
kimianya bila tidak; lalu bab ini mengelompokkan batuan itu ke dalam seri, yaitu keluarga batuan yang
berkerabat melalui diferensiasi, dan mengaitkan setiap seri dengan suatu tatanan tektonik, beserta
diagram yang dipakai ahli petrologi untuk membaca tatanan dari sebuah analisis, dan batas-batasnya.
% ledger: vp21:series

\begin{recall}
\textbf{Dari volume sekolah menengah atas}: basal, andesit, riolit, gabro, dan granit;
batuan felsik dan mafik; tekstur vulkanik dan plutonik. \textbf{Dari volume Tahun
1}: punggungan tengah samudra, zona subduksi dan busur vulkanik, titik panas dan provinsi
beku besar, rift benua. \textbf{Dari \cref{ch:b2:elements-in-the-earth}}:
MORB dan OIB, \omterm{def:b2:elements-in-the-earth:normalised}{diagram ternormalisasi} \omterm{def:b2:elements-in-the-earth:bse}{mantel primitif} (diagram laba-laba), \omterm{def:b2:elements-in-the-earth:partition}{unsur kompatibel} dan
inkompatibel. \textbf{Dari \cref{ch:b2:isotope-geochemistry}}: \omterm{def:b2:isotope-geochemistry:reservoirs}{mantel
terkuras} dan reservoir isotop yang diperkaya. \textbf{Dari
\cref{ch:b2:magmatic-differentiation}}: \omterm{def:b2:magmatic-differentiation:crystallisation}{kristalisasi fraksional}, \omterm{def:b2:magmatic-differentiation:lld}{garis turunan
cairan}, \omterm{def:b2:magmatic-differentiation:harker}{diagram Harker}, dan hukum Rayleigh.
\end{recall}

\section{Menamai batuan beku}\label{sec:b2:magmatic-series:naming}

Batuan berbutir kasar memperlihatkan mineralnya, dan dinamai berdasarkan mineral itu. Lava sering kali terlalu
halus butirannya atau terlalu banyak kacanya untuk itu, sehingga dinamai berdasarkan analisis kimianya. Kedua
skema itu adalah konvensi internasional: keduanya memerikan, dan tidak menyiratkan apa pun tentang bagaimana
batuan itu terbentuk.
% ledger: bgs:tas-use

\begin{definition}[Komposisi modal; diagram QAPF]\label{def:b2:magmatic-series:qapf}
Yang disebut \emph{komposisi modal}\index{komposisi modal} suatu batuan adalah proporsi menurut
volume setiap mineralnya. Yang disebut \emph{diagram QAPF}\index{diagram QAPF} adalah diagram yang menamai
batuan beku berbutir kasar dengan mineral mafik kurang dari \qty{90}{\percent} berdasarkan
empat komponen modal yang dihitung ulang sehingga berjumlah 100: kuarsa (Q), feldspar alkali
(A), plagioklas (P), dan feldspatoid (F). Karena kuarsa dan feldspatoid tidak
mengkristal bersama, setiap batuan terletak di salah satu dari dua segitiga, QAP atau FAP, yang bersambung
pada sisi bersamanya, AP.
\end{definition}
% ledger: bgs:qapf

\begin{omfigure}
\centering
\begin{tikzpicture}[font=\scriptsize, x=1cm, y=1cm]
\def\L{6.4}
\pgfmathsetmacro{\H}{\L*0.8660254}
% upper point (q, p): q = Q fraction, p = P/(A+P)
\newcommand{\qapfup}[2]{({(#1)*\L/2 + (#2)*(1-(#1))*\L}, {(#1)*\H})}
\newcommand{\dn}[2]{({(#1)*\L/2 + (#2)*(1-(#1))*\L}, {-(#1)*\H})}
\draw[thick] (0,0) -- (\L/2,\H) -- (\L,0) -- (\L/2,-\H) -- cycle;
\draw (0,0) -- (\L,0);
\foreach \q in {0.9,0.6,0.2,0.05} \draw \qapfup{\q}{0} -- \qapfup{\q}{1};
\foreach \q in {0.1,0.6} \draw \dn{\q}{0} -- \dn{\q}{1};
\foreach \p in {0.1,0.65,0.9} \draw \qapfup{0.6}{\p} -- \qapfup{0}{\p};
\draw[dashed] \qapfup{0.6}{0.35} -- \qapfup{0.2}{0.35};
\draw \qapfup{0.2}{0.35} -- \qapfup{0}{0.35};
\foreach \p in {0.1,0.35,0.65,0.9} \draw \dn{0}{\p} -- \dn{0.1}{\p};
\foreach \p in {0.1,0.5,0.9} \draw \dn{0.1}{\p} -- \dn{0.6}{\p};
\node[anchor=south] at (\L/2,\H) {\normalsize Q};
\node[anchor=east] at (0,0) {\normalsize A};
\node[anchor=west] at (\L,0) {\normalsize P};
\node[anchor=north] at (\L/2,-\H) {\normalsize F};
\node at \qapfup{0.94}{0.5} {1a};
\node at \qapfup{0.75}{0.5} {1b};
\node[font=\tiny] at \qapfup{0.4}{0.05} {2};
\node at \qapfup{0.4}{0.225} {3a};
\node at \qapfup{0.4}{0.5} {3b};
\node at \qapfup{0.4}{0.775} {4};
\node[font=\tiny] at \qapfup{0.4}{0.95} {5};
\foreach \p/\n in {0.225/7*,0.5/8*,0.775/9*} \node at \qapfup{0.125}{\p} {\n};
\node[font=\tiny] at \qapfup{0.125}{0.05} {6*};
\node[font=\tiny] at \qapfup{0.125}{0.95} {10*};
\foreach \p/\n in {0.225/7,0.5/8,0.775/9} \node at \qapfup{0.025}{\p} {\n};
\foreach \p/\n in {0.225/7',0.5/8',0.775/9'} \node at \dn{0.05}{\p} {\n};
\node[font=\tiny] at \dn{0.3}{0.05} {11};
\node at \dn{0.35}{0.3} {12};
\node at \dn{0.35}{0.7} {13};
\node[font=\tiny] at \dn{0.3}{0.95} {14};
\node at \dn{0.75}{0.5} {15};
\node[anchor=east] at \qapfup{0.9}{0} {90};
\node[anchor=east] at \qapfup{0.6}{0} {60};
\node[anchor=east] at \qapfup{0.2}{0} {20};
\node[anchor=east] at \dn{0.1}{0} {10};
\node[anchor=east] at \dn{0.6}{0} {60};
\foreach \p/\n in {0.1/10,0.35/35,0.65/65,0.9/90} \node[anchor=south, fill=white, inner sep=0.5pt] at \qapfup{0.6}{\p} {\tiny\n};
\node[anchor=west, align=left, font=\scriptsize] at (\L+0.9,1.6) {Segitiga atas: \% Q\\dalam Q + A + P;\\
  garis-garis\\$100\,P/(A+P)$ tetap\\pada 10, 35, 65, 90};
\node[anchor=west, align=left, font=\scriptsize] at (\L+0.9,-1.7) {Segitiga bawah: \% F\\dalam F + A + P;\\
  $100\,P/(A+P)$\\pada 10, 50, 90};
\end{tikzpicture}
\par\smallskip{\small Klasifikasi QAPF batuan beku berbutir kasar. 1a
kuarsolit, 1b granitoid kaya kuarsa, 2 granit feldspar alkali, 3a sienogranit dan
3b monzogranit (bersama-sama disebut granit), 4 granodiorit, 5 tonalit; 6*--10* varietas
mengandung kuarsa dari 6--10: 6 sienit feldspar alkali, 7 sienit, 8 monzonit, 9 monzodiorit
atau monzogabro, 10 diorit, gabro, atau anortosit; 6'--10' varietasnya yang mengandung
foid; 11 sienit foid, 12 monzosienit foid, 13 monzodiorit foid atau
monzogabro foid, 14 diorit foid atau gabro foid, 15 foidolit. Medan-medan tipis 6, 6', 10, dan 10'
di sepanjang tepi tidak diberi nomor. Menurut skema British Geological Survey (berdasarkan
Streckeisen 1976).}
\end{omfigure}
% ledger: bgs:qapf

Medan 10 memuat tiga nama. \omterm{def:b2:magmatic-series:qapf}{Diagram QAPF} tidak memisahkan diorit, gabro, dan
anortosit, juga tidak memisahkan batuan gabroik satu sama lain; diagram-diagram lain, yang berdasarkan
plagioklas dan mineral mafik, melakukannya.

\begin{definition}[Diagram alkali total--silika]\label{def:b2:magmatic-series:tas}
Yang disebut \emph{diagram alkali total--silika}\index{diagram alkali total--silika} (diagram
TAS) adalah diagram yang menamai batuan vulkanik yang modusnya tidak dapat diukur berdasarkan persentase
massa \ce{Na2O + K2O} terhadap \ce{SiO2}, setelah analisisnya dihitung
ulang menjadi \qty{100}{\percent} tanpa air dan karbon dioksida. Diagram ini dibagi
menjadi lima belas medan oleh garis-garis lurus yang titik potongnya berkoordinat tetap
(gambar sesudah paragraf berikut).
\end{definition}
% ledger: bgs:tas, bgs:tas-use

Diagram ini memiliki tiga kelompok medan. Di sepanjang bagian bawah, pada kadar alkali rendah, terletak
nama-nama seri yang umum, pikrobasal dan basal sampai riolit, yang dipisahkan pada silika 45, 52,
57, dan \qty{63}{\percent} dan, untuk riolit, oleh sebuah garis miring. Di atasnya, medan S dan trakit memuat batuan yang lebih kaya akan
alkali pada silika yang sama; medan U, fonolit, dan foidit memuat batuan yang begitu kaya akan
alkali, dan begitu miskin silika, sehingga mengkristalkan feldspatoid. Beberapa medan
dibagi lagi dengan kriteria kedua: trakibasal disebut hawaiit jika $\ce{Na2O} - 2
\ge \ce{K2O}$, dan trakibasal kalium jika tidak; tefrit dengan olivin lebih dari
\qty{10}{\percent} dalam normanya (lihat di bawah) adalah basanit; riolit atau trakit yang
rasio molekuler $(\ce{Na2O + K2O})/\ce{Al2O3}$-nya melebihi 1 disebut peralkali. Skema ini
menuntut analisis dengan air kurang dari \qty{2}{\percent} dan karbon dioksida kurang dari \qty{0.5}{\percent}:
alterasi menambahkan air dan memindahkan alkali, dan \omterm{def:b2:magmatic-differentiation:cumulate}{kumulat} bukanlah
cairan.
% ledger: bgs:tas, bgs:tas-use

\begin{omfigure}
\centering
\begin{tikzpicture}
\begin{axis}[omrecord, width=0.86\linewidth, height=8.6cm,
  xmin=37, xmax=80, ymin=0, ymax=16, xtick={37,41,...,77}, ytick={0,2,...,16},
  xlabel={\ce{SiO2} (\unit{\percent})}, ylabel={\ce{Na2O + K2O} (\unit{\percent})},
  legend style={font=\scriptsize, at={(0.98,0.98)}, anchor=north east},
  tick label style={font=\scriptsize}, label style={font=\small}]
\addplot[only marks, mark=*, mark size=0.55pt, blue!70!black, opacity=0.6] table[x=SiO2, y=alk] {figdata/out/bachelor-2-series-cath-th.dat};
\addlegendentry{Islandia}
\addplot[only marks, mark=*, mark size=0.55pt, red!75!black, opacity=0.6] table[x=SiO2, y=alk] {figdata/out/bachelor-2-series-cath-ca.dat};
\addlegendentry{Cascade}
\foreach \f in {Pc,B,O1,O2,O3,R,S1,S2,S3,T,U1,U2,U3,Ph,F}
  \addplot[black, thin, forget plot] table[x=SiO2, y=alk] {figdata/out/bachelor-2-series-tas-\f.dat};
\node[font=\tiny] at (axis cs:43,1.6) {Pc};
\node[font=\tiny, fill=white, inner sep=0.5pt] at (axis cs:50.5,0.8) {B};
\node[font=\tiny] at (axis cs:54.5,2.2) {O1};
\node[font=\tiny] at (axis cs:60,2.2) {O2};
\node[font=\tiny] at (axis cs:66.5,2.2) {O3};
\node[font=\tiny] at (axis cs:74.5,11) {R};
\node[font=\tiny, fill=white, inner sep=0.5pt] at (axis cs:49,5.8) {S1};
\node[font=\tiny] at (axis cs:52.5,7.4) {S2};
\node[font=\tiny] at (axis cs:57.5,8.9) {S3};
\node[font=\tiny] at (axis cs:64.5,12) {T};
\node[font=\tiny] at (axis cs:44,6.6) {U1};
\node[font=\tiny] at (axis cs:48.7,9.4) {U2};
\node[font=\tiny] at (axis cs:52.5,11.8) {U3};
\node[font=\tiny] at (axis cs:57,14.6) {Ph};
\node[font=\tiny] at (axis cs:39,10) {F};
\end{axis}
\end{tikzpicture}
\par\smallskip{\small Diagram TAS (menurut Le Bas dan kawan-kawan, 1986, sebagaimana disajikan oleh
British Geological Survey) dengan 456 lava dan batuan intrusif dari Islandia (biru) dan
580 dari Cascade (merah), oksida sebagaimana dianalisis. Pc pikrobasal, B basal, O1 andesit
basaltik, O2 andesit, O3 dasit, R riolit; S1 trakibasal, S2 trakiandesit
basaltik, S3 trakiandesit, T trakit atau trakidasit; U1 tefrit atau
basanit, U2 fonotefrit, U3 tefrifonolit, Ph fonolit, F foidit.}
\end{omfigure}
% ledger: bgs:tas, geoplotr:tas, data:b4-cath


\begin{definition}[Kejenuhan silika]\label{def:b2:magmatic-series:saturation}
\omterm{def:b2:magma-emplacement:magma}{Magma} atau batuan disebut \emph{lewat jenuh silika}\index{kejenuhan silika} jika mengandung
silika lebih banyak daripada yang dapat digabungkan dengan oksida-oksida lainnya dalam mineral silikat, sehingga
kuarsa (atau mineral silika lain) muncul; \emph{kurang jenuh silika} jika silikanya terlalu
sedikit untuk mengubah semua alkalinya menjadi feldspar, sehingga feldspatoid (nefelin,
leusit) yang muncul, tidak pernah bersama kuarsa; dan \emph{jenuh silika}
di antara keduanya.
\end{definition}

Untuk lava yang mineralnya tidak dapat dihitung, kejenuhan dinilai berdasarkan \emph{norma}, yaitu
perhitungan baku yang membagikan oksida hasil analisis ke dalam mineral-mineral ideal
dengan urutan tetap, alkali ke dalam feldspar lebih dahulu, lalu melaporkan kelebihan silika sebagai
kuarsa atau kekurangannya sebagai feldspatoid. Norma yang lazim, norma CIPW, memisahkan
dua jenis basal: basal dengan nefelin dalam normanya adalah basal alkali; basal
tanpa nefelin adalah basal subalkali, dan di antaranya basal toleiitik mengandung
hipersten (ortopiroksen) dalam normanya dan, sering kali, piroksen rendah kalsium dalam modusnya.
\Cref{exo:b2:magmatic-series:11} menjalankan versi sederhananya.
% ledger: bgs:tas-use, bgs:series-terms

\section{Seri magmatik}\label{sec:b2:magmatic-series:series}

Lava dari satu gunung api, atau dari satu provinsi, jarang sama, tetapi lava-lava itu berkerabat: lava itu
berasal dari \omterm{def:b2:magma-emplacement:magma}{magma} induk yang serupa melalui diferensiasi pada
\cref{ch:b2:magmatic-differentiation}, dan bersama-sama menelusuri satu garis turunan.

\begin{definition}[Seri magmatik]\label{def:b2:magmatic-series:series}
Yang disebut \emph{seri magmatik}\index{seri magmatik} adalah rangkaian batuan beku dari satu
siklus atau provinsi magmatik yang berkerabat melalui diferensiasi dari \omterm{def:b2:magma-emplacement:magma}{magma} induk yang serupa dan
memiliki ciri kimia dan mineralogi yang khas. Tiga seri yang lazim
dibedakan. \emph{Seri toleiitik}\index{seri toleiitik} berawal dari
basal yang jenuh atau sedikit \omterm{def:b2:magmatic-series:saturation}{lewat jenuh silika}, dengan hipersten dalam normanya,
dan menunjukkan pengayaan besi yang nyata pada awal diferensiasinya.
\emph{Seri kalk-alkali}\index{seri kalk-alkali} berkisar dari basal melalui
andesit sampai dasit dan riolit, dengan sedikit atau tanpa pengayaan besi, dan dengan
mineral hidrus (hornblende, biotit) pada anggota-anggotanya yang telah berevolusi. \emph{Seri
alkali}\index{seri alkali} berawal dari basal bernorma nefelin dan berevolusi
melalui medan S dan U pada diagram TAS ke arah trakit dan fonolit.
\end{definition}
% ledger: bgs:series-terms, bgs:tas-use

Nama-nama ini memerikan rangkaian batuan, bukan batuan tunggal: sebuah basal yang terpisah dapat termasuk ke salah satu
dari ketiganya, dan nama-nama yang sama telah dipakai dengan arti yang berbeda-beda, sehingga ahli petrologi
harus menyatakan apa yang ia maksud dengan nama itu. Perbedaan antara \omterm{def:b2:magmatic-series:series}{seri toleiitik} dan
\omterm{def:b2:magmatic-series:series}{seri kalk-alkali} bertumpu pada besi, dan diagram yang paling baik memperlihatkannya memplot tiga
kelompok oksida.
% ledger: bgs:series-terms

\begin{definition}[Diagram AFM]\label{def:b2:magmatic-series:afm}
Yang disebut \emph{diagram AFM}\index{diagram AFM} adalah diagram segitiga yang sudut-sudutnya
$A = \ce{Na2O + K2O}$, $F = \ce{FeO_T}$ (seluruh besi sebagai \ce{FeO}), dan $M = \ce{MgO}$, dalam
persen massa yang dihitung ulang sehingga $A + F + M = 100$. Sebuah garis yang ditarik padanya memisahkan
medan toleiitik, di sisi F, dari medan kalk-alkali.
\end{definition}
% ledger: vp21:series, geoplotr:afm

Basal yang terdiferensiasi kehilangan magnesium lebih cepat daripada unsur lainnya, karena olivin dan
piroksen lebih kaya akan magnesium daripada cairan, dan mempertahankan alkalinya, yang bersifat
inkompatibel: garis turunannya bergerak menjauhi M, ke arah A. Nasib
besi menentukan serinya. Pada \omterm{def:b2:magma-emplacement:magma}{magma} toleiitik, besi hanya keluar bersama silikat
fero yang kaya magnesium, lebih lambat daripada magnesium, sehingga rasio
$F/M$ naik dan tren itu mula-mula naik ke arah F sebelum berbelok ke A:
\emph{tren Fenner}. Pada \omterm{def:b2:magma-emplacement:magma}{magma} kalk-alkali, yang lebih teroksidasi, besi juga keluar dalam
magnetit, yang mengandung besi feri, dan $F/M$ tetap hampir konstan sementara cairan
bergerak lurus ke arah A: \emph{tren Bowen}. Gambar di bawah memperlihatkan keduanya dengan
batuan nyata, toleiit Islandia dan lava kalk-alkali Cascade. Garis
pemisahnya, yang ditarik secara visual pada 1971, memisahkan keduanya dengan baik di bagian basaltik; pada
kadar alkali tinggi kedua tren itu menyatu dan garis itu tidak banyak berarti.
% ledger: vp21:series, geoplotr:afm, data:b4-cath

\begin{omfigure}
\centering
\begin{tikzpicture}
\begin{axis}[name=afm, width=0.5\linewidth, height=0.5\linewidth, axis equal image,
  hide axis, xmin=-0.04, xmax=1.04, ymin=-0.06, ymax=0.92, clip=false]
\addplot[only marks, mark=*, mark size=0.5pt, blue!70!black, opacity=0.6] table[x=afm_x, y=afm_y] {figdata/out/bachelor-2-series-cath-th.dat};
\addplot[only marks, mark=*, mark size=0.5pt, red!75!black, opacity=0.6] table[x=afm_x, y=afm_y] {figdata/out/bachelor-2-series-cath-ca.dat};
\addplot[very thick, black] table[x=x, y=y] {figdata/out/bachelor-2-series-afm-line.dat};
\draw[thick] (axis cs:0,0) -- (axis cs:1,0) -- (axis cs:0.5,0.8660254) -- cycle;
\node[anchor=south] at (axis cs:0.5,0.866) {F};
\node[anchor=north east] at (axis cs:0,0) {A};
\node[anchor=north west] at (axis cs:1,0) {M};
\node[font=\scriptsize, fill=white, inner sep=1pt] at (axis cs:0.52,0.62) {toleiitik};
\node[font=\scriptsize, fill=white, inner sep=1pt] at (axis cs:0.33,0.12) {kalk-alkali};
\end{axis}
\begin{axis}[omrecord, at={(afm.east)}, anchor=west, xshift=1.3cm,
  width=0.46\linewidth, height=0.46\linewidth,
  xmin=-2.5, xmax=2.2, ymin=-4.5, ymax=1.5,
  xlabel={$\ln(A/F)$}, ylabel={$\ln(M/F)$},
  legend style={font=\scriptsize, at={(0.98,0.98)}, anchor=north east},
  tick label style={font=\scriptsize}, label style={font=\small}]
\addplot[only marks, mark=*, mark size=0.5pt, blue!70!black, opacity=0.6] table[x=lnAF, y=lnMF] {figdata/out/bachelor-2-series-lr-th.dat};
\addlegendentry{Islandia}
\addplot[only marks, mark=*, mark size=0.5pt, red!75!black, opacity=0.6] table[x=lnAF, y=lnMF] {figdata/out/bachelor-2-series-lr-ca.dat};
\addlegendentry{Cascade}
\addplot[thick, blue!50!black, forget plot] table[x=x, y=y_th] {figdata/out/bachelor-2-series-lr-fits.dat};
\addplot[thick, red!50!black, forget plot] table[x=x, y=y_ca] {figdata/out/bachelor-2-series-lr-fits.dat};
\end{axis}
\end{tikzpicture}
\par\smallskip{\small Kiri: \omterm{def:b2:magmatic-series:afm}{diagram AFM} batuan Islandia (biru) dan Cascade (merah)
pada gambar TAS, dengan garis pemisah Irvine dan Baragar (1971). Kanan:
batuan yang sama dalam koordinat log-rasio, tempat setiap seri menjadi pita yang hampir
lurus; garis-garisnya adalah regresi ortogonal, dengan kemiringan $-1.10$ (Islandia) dan $-0.72$
(Cascade), yang bertemu di $(-3.84, 2.12)$ seperti yang ditemukan Vermeesch dan Pease (2021).}
\end{omfigure}
% ledger: geoplotr:afm, data:b4-cath, vp21:model

Dari batuan Islandia yang berada dalam jangkauan garis pemisah, 351 jatuh di sisi
toleiitik dan 23 di sisi lain; dari batuan Cascade, 417 di sisi kalk-alkali
dan 57 di sisi lain. Kurva adalah batas yang canggung, dan sebuah model sederhana menjelaskan
bentuk tren-tren itu. Misalkan \omterm{def:b2:magma-emplacement:magma}{magma} kehilangan alkali, besi, dan magnesiumnya
dengan laju yang sebanding dengan jumlah yang dikandungnya, $\mathrm dX/\mathrm dt = -\lambda_X X$.

\begin{proposition}[Tren log-rasio]\label{prop:b2:magmatic-series:logratio}
Jika $A$, $F$, dan $M$ meluruh secara eksponensial dengan konstanta $\lambda_A$, $\lambda_F$, dan
$\lambda_M$, \omterm{def:b2:magma-emplacement:magma}{magma} bergerak menyusuri garis lurus pada bidang $\ln(A/F)$ dan
$\ln(M/F)$, dengan kemiringan $(\lambda_F - \lambda_M)/(\lambda_F - \lambda_A)$, sementara \omterm{def:b2:magma-emplacement:magma}{magma} itu bergerak
menyusuri kurva pada segitiga AFM.
\end{proposition}

\begin{proof}
Dengan $X = X_0 e^{-\lambda_X t}$, $\ln(A/F) = \ln(A_0/F_0) + (\lambda_F - \lambda_A)t$
dan $\ln(M/F) = \ln(M_0/F_0) + (\lambda_F - \lambda_M)t$: keduanya linear dalam $t$, sehingga
yang satu linear terhadap yang lain, dengan kemiringan yang dinyatakan tadi. Rasio $A/F$ dan $M/F$ menentukan
titik pada segitiga, yang koordinatnya $A/(A + F + M)$, \dots\ tidak linear dalam $t$.
\end{proof}

Kemiringan batuan Islandia yang lebih curam berarti $\lambda_F$ yang lebih kecil: besi meninggalkan
\omterm{def:b2:magma-emplacement:magma}{magma} toleiitik lebih lambat (\cref{exo:b2:magmatic-series:10}). Vermeesch dan Pease menunjukkan bahwa kedua
pita pada gambar di atas bertemu di satu titik, yang mereka tafsirkan sebagai sumber bersama
yang hipotetis.
% ledger: vp21:model

\omterm{def:b2:magmatic-series:series}{Seri toleiitik} dan \omterm{def:b2:magmatic-series:series}{seri kalk-alkali} juga berbeda dalam kadar kalium oksidanya. Diagram
\ce{K2O}--\ce{SiO2} membagi basal, andesit, dasit, dan riolit dengan
silika lebih dari \qty{48}{\percent} menjadi tipe K-rendah, K-sedang, dan K-tinggi dengan dua
garis lurus (gambar di bawah). Sebagian besar kedua rangkaian di sini bertipe K-sedang:
diagram itu menyatakan berapa banyak kalium oksida yang dibawa suatu seri, bukan seri apa itu.
% ledger: bgs:k2o-sio2, data:b4-cath

\begin{omfigure}
\centering
\begin{tikzpicture}
\begin{axis}[omrecord, width=0.78\linewidth, height=6.2cm,
  xmin=45, xmax=79, ymin=0, ymax=5,
  xlabel={\ce{SiO2} (\unit{\percent})}, ylabel={\ce{K2O} (\unit{\percent})},
  legend style={font=\scriptsize, at={(0.02,0.98)}, anchor=north west},
  tick label style={font=\scriptsize}, label style={font=\small}]
\addplot[only marks, mark=*, mark size=0.55pt, blue!70!black, opacity=0.6] table[x=SiO2, y=K2O] {figdata/out/bachelor-2-series-cath-th.dat};
\addlegendentry{Islandia}
\addplot[only marks, mark=*, mark size=0.55pt, red!75!black, opacity=0.6] table[x=SiO2, y=K2O] {figdata/out/bachelor-2-series-cath-ca.dat};
\addlegendentry{Cascade}
\addplot[thick, black, forget plot] table[x=x, y=y_low] {figdata/out/bachelor-2-series-k2o-lines.dat};
\addplot[thick, black, forget plot] table[x=x, y=y_high] {figdata/out/bachelor-2-series-k2o-lines.dat};
\draw[thick] (axis cs:48,0) -- (axis cs:48,0.3);
\node[font=\scriptsize, fill=white, inner sep=1pt] at (axis cs:60,3.7) {K-tinggi};
\node[font=\scriptsize, fill=white, inner sep=1pt] at (axis cs:76.5,2.2) {K-sedang};
\node[font=\scriptsize, fill=white, inner sep=1pt] at (axis cs:60,0.35) {K-rendah};
\end{axis}
\end{tikzpicture}
\par\smallskip{\small Kalium oksida terhadap silika untuk batuan Islandia dan Cascade, dengan
garis melalui (48, 0.3)--(68, 1.2) dan (48, 1.2)--(68, 3.1) yang memisahkan tipe K-rendah,
K-sedang, dan K-tinggi (menurut Peccerillo dan Taylor, 1976, sebagaimana diadopsi oleh
skema internasional), yang di sini diperpanjang melewati 68.}
\end{omfigure}
% ledger: bgs:k2o-sio2, data:b4-cath

Mengapa tingkat oksidasinya berbeda? \omterm{def:b2:magma-emplacement:magma}{Magma} kalk-alkali terbentuk di atas zona subduksi, tempat
slab yang melepaskan air memasok spesies teroksidasi ke baji mantel; \omterm{def:b2:magma-emplacement:magma}{magma} itu
mengkristalkan magnetit, dan anggota-anggotanya yang telah berevolusi mengkristalkan hornblende atau biotit. \omterm{def:b2:magma-emplacement:magma}{Magma}
toleiitik terbentuk di punggungan, tempat mantel meleleh melalui dekompresi
(\cref{ch:b2:mantle-melting}), dan lebih tereduksi.
% ledger: vp21:series

\section{Seri dan tatanan}\label{sec:b2:magmatic-series:settings}

Setiap tatanan tektonik lempeng menghasilkan \omterm{def:b2:magma-emplacement:magma}{magmanya} sendiri, karena masing-masing melelehkan
mantel dengan caranya sendiri: melalui dekompresi di bawah punggungan dan titik panas, melalui penambahan air
di atas zona subduksi (\cref{ch:b2:mantle-melting}).

\emph{Punggungan tengah samudra} meletuskan basal toleiitik, yaitu MORB pada
\cref{ch:b2:elements-in-the-earth}, yang miskin \omterm{def:b2:elements-in-the-earth:partition}{unsur inkompatibel} dan membawa
tanda isotop terkuras dari sumbernya. Islandia, punggungan di atas titik panas, menambahkan
toleiitnya sendiri ke dalam daftar itu, dengan rangkaian lengkap sampai riolit.
% ledger: vp21:series

\emph{Pulau samudra} tumbuh secara bertahap. Gunung api Hawaii meletuskan lava alkali ketika
masih muda, lalu membangun perisainya dari basal toleiitik, yang diperkirakan 80 sampai
\qty{95}{\percent} volume akhirnya, dalam 1 sampai 2 juta tahun pertamanya; ketika
gunung api itu hanyut menjauhi titik panas, lavanya kembali menjadi alkali, dan berakhir dengan ekor panjang
erupsi alkali yang jarang. Pelelehan yang lebih sedikit menghasilkan \omterm{def:b2:magma-emplacement:magma}{magma} yang lebih alkali: alkali, yang
inkompatibel, lebih sedikit diencerkan dalam \omterm{def:b2:mantle-melting:fraction}{fraksi lelehan} yang kecil. Secara isotopik, \omterm{def:b2:elements-in-the-earth:morb-oib}{basal
pulau samudra} mencuplik reservoir yang diperkaya pada \cref{ch:b2:isotope-geochemistry}.
% ledger: usgs14:hawaii-stages

\emph{Busur vulkanik} meletuskan \omterm{def:b2:magmatic-series:series}{seri kalk-alkali}. \omterm{def:b2:elements-in-the-earth:trace}{Unsur jejaknya} membawa
tanda yang khas (gambar di bawah): bila dinormalkan terhadap \omterm{def:b2:elements-in-the-earth:bse}{mantel primitif}, basal
busur diperkaya dalam Rb, Ba, Th, U, Pb, dan Sr, unsur-unsur yang ditambahkan slab ke
baji, dan terkuras dalam niobium dan tantalum dibandingkan dengan unsur tetangganya. Basal busur
median dalam himpunan data itu memiliki $\mathrm{Nb}/\mathrm{La} = 0.49$, dibandingkan dengan 0.97 untuk MORB
dan 1.23 untuk \omterm{def:b2:elements-in-the-earth:morb-oib}{basal pulau samudra}.
% ledger: data:b4-basalt-settings

\begin{omfigure}
\centering
\begin{tikzpicture}
\begin{axis}[omrecord, width=0.86\linewidth, height=6.4cm, ymode=log,
  xmin=0.5, xmax=22.5, ymin=1, ymax=300,
  xtick={1,...,22}, xticklabels={Rb,Ba,Th,U,Nb,Ta,La,Ce,Pb,Sr,Nd,Zr,Hf,Sm,Eu,Ti,Gd,Dy,Y,Er,Yb,Lu},
  ylabel={cuplikan / mantel primitif},
  legend style={font=\scriptsize, at={(0.98,0.98)}, anchor=north east},
  tick label style={font=\scriptsize}, label style={font=\small}]
\addplot[thick, mark=square*, mark size=1.4pt, orange!85!black] table[x=idx, y=oib] {figdata/out/bachelor-2-series-spider.dat};
\addlegendentry{basal pulau samudra}
\addplot[thick, mark=*, mark size=1.4pt, blue!70!black] table[x=idx, y=morb] {figdata/out/bachelor-2-series-spider.dat};
\addlegendentry{MORB}
\addplot[thick, mark=triangle*, mark size=1.7pt, red!75!black] table[x=idx, y=iab] {figdata/out/bachelor-2-series-spider.dat};
\addlegendentry{basal busur kepulauan}
\end{axis}
\end{tikzpicture}
\par\smallskip{\small Komposisi median ternormalisasi \omterm{def:b2:elements-in-the-earth:bse}{mantel primitif} dari 241 MORB, 259
\omterm{def:b2:elements-in-the-earth:morb-oib}{basal pulau samudra}, dan 256 basal busur kepulauan (kompilasi Vermeesch, 2006;
\omterm{def:b2:elements-in-the-earth:bse}{mantel primitif} menurut McDonough dan Sun, 1995). Basal busur memperlihatkan lembah pada Nb
dan Ta serta puncak pada Ba, Pb, dan Sr; \omterm{def:b2:elements-in-the-earth:morb-oib}{basal pulau samudra} diperkaya dalam semua
unsur yang paling inkompatibel.}
\end{omfigure}
% ledger: data:b4-basalt-settings, ms95:pm-trace

Di dalam benua, \omterm{def:b2:magma-emplacement:magma}{magma} basaltik bertemu dengan kerak tebal yang dapat dipanaskan dan dilelehkannya,
dan hasilnya dapat berupa vulkanisme dengan dua modus.

\begin{definition}[Vulkanisme bimodal]\label{def:b2:magmatic-series:bimodal}
Yang disebut \emph{vulkanisme bimodal}\index{vulkanisme bimodal} adalah erupsi, dalam satu provinsi,
dua kelompok \omterm{def:b2:magma-emplacement:magma}{magma}, basaltik dan riolitik, dengan sedikit komposisi menengah
di antara keduanya.
\end{definition}

Vulkanisme di sepanjang jejak titik panas Yellowstone bersifat bimodal: di Dataran
Sungai Snake bagian timur, riolit yang lebih tua dengan silika sekitar \qty{75}{\percent} tertutup oleh
basal yang lebih muda dengan silika sekitar \qty{50}{\percent}. Celah di antara keduanya adalah petunjuknya:
diferensiasi basal saja akan meninggalkan lava menengah, dan celah itu
menunjuk pada sumber kedua bagi riolit, misalnya kerak yang dipanaskan basal.
% ledger: yvo23:bimodal

\emph{Provinsi beku besar}, yang telah dijumpai dalam volume Tahun 1, meletuskan basal dalam volume
yang tidak tertandingi oleh tatanan lain; serinya dibaca dengan diagram yang sama.

\section{Diagram diskriminasi}\label{sec:b2:magmatic-series:discrimination}

Pola-pola pada subbab sebelumnya mengisyaratkan bahwa tatanan dapat dibaca langsung dari analisis,
dan ahli petrologi telah menggambar banyak diagram untuk itu.

\begin{definition}[Diagram diskriminasi]\label{def:b2:magmatic-series:discrimination}
Yang disebut \emph{diagram diskriminasi}\index{diagram diskriminasi} adalah diagram
variabel kimia, biasanya \omterm{def:b2:elements-in-the-earth:trace}{unsur jejak} yang tidak mudah berpindah, yang dibagi menjadi medan-medan yang ditarik
di sekeliling komposisi batuan dari tatanan yang diketahui, dan dipakai untuk menetapkan tatanan bagi
batuan yang asal-usulnya tidak diketahui.
\end{definition}

Contoh klasiknya adalah segitiga Ti--Zr--Y dari Pearce dan Cann (1973), yang dibuat untuk basal. Diagram itu
memakai unsur yang sedikit dipindahkan oleh alterasi, sehingga dapat diterapkan pada lava tua
yang termetamorfosis dan tatanannya telah hilang. Gambar di bawah mengujinya pada 538 basal dalam
kompilasi itu yang memiliki ketiga unsur tersebut. Sebagian besar \omterm{def:b2:elements-in-the-earth:morb-oib}{basal pulau samudra} jatuh di
medan dalam-lempeng D, 152 dari 173, dan sebagian besar MORB di medan B, 142 dari 164; tetapi medan B
juga ditempati basal busur, dan 101 dari 201 basal busur juga jatuh di sana.
% ledger: geoplotr:tizry, data:b4-basalt-settings

\begin{omfigure}
\centering
\begin{tikzpicture}
\begin{axis}[width=0.62\linewidth, height=0.55\linewidth, axis equal image,
  hide axis, xmin=-0.04, xmax=1.04, ymin=-0.06, ymax=0.92, clip=false,
  legend style={font=\scriptsize, at={(1.0,0.95)}, anchor=north west}]
\addplot[only marks, mark=*, mark size=0.7pt, blue!70!black, opacity=0.7] table[x=x, y=y] {figdata/out/bachelor-2-series-tizry-morb.dat};
\addlegendentry{MORB}
\addplot[only marks, mark=square*, mark size=0.7pt, orange!85!black, opacity=0.7] table[x=x, y=y] {figdata/out/bachelor-2-series-tizry-oib.dat};
\addlegendentry{basal pulau samudra}
\addplot[only marks, mark=triangle*, mark size=0.9pt, red!75!black, opacity=0.7] table[x=x, y=y] {figdata/out/bachelor-2-series-tizry-iab.dat};
\addlegendentry{basal busur kepulauan}
\foreach \f in {A,B,C,D}
  \addplot[thick, black, forget plot] table[x=x, y=y] {figdata/out/bachelor-2-series-tizry-\f.dat};
\draw[thick] (axis cs:0,0) -- (axis cs:1,0) -- (axis cs:0.5,0.8660254) -- cycle;
\node[anchor=south] at (axis cs:0.5,0.866) {Ti/100};
\node[anchor=north east] at (axis cs:0,0) {Zr};
\node[anchor=north west] at (axis cs:1,0) {Y $\times$ 3};
\draw[thin] (axis cs:0.33,0.36) -- (axis cs:0.22,0.46) node[anchor=south east, font=\small\bfseries, inner sep=1pt] {D};
\draw[thin] (axis cs:0.36,0.12) -- (axis cs:0.24,0.06) node[anchor=north east, font=\small\bfseries, inner sep=1pt] {C};
\draw[thin] (axis cs:0.56,0.37) -- (axis cs:0.76,0.50) node[anchor=west, font=\small\bfseries, inner sep=1pt] {A};
\draw[thin] (axis cs:0.53,0.21) -- (axis cs:0.72,0.13) node[anchor=north west, font=\small\bfseries, inner sep=1pt] {B};
\end{axis}
\end{tikzpicture}
\par\smallskip{\small Diagram Ti--Zr--Y dari Pearce dan Cann (1973) dengan basal-basal
yang tatanannya diketahui. A toleiit busur kepulauan; B MORB, toleiit busur kepulauan, dan
basal kalk-alkali; C basal kalk-alkali; D basal dalam-lempeng. Ti dari
\ce{TiO2}; ketiganya dalam ppm.}
\end{omfigure}
% ledger: geoplotr:tizry, data:b4-basalt-settings

\begin{method}[Memakai diagram diskriminasi]\label{met:b2:magmatic-series:discrimination}
(1) Periksalah bahwa batuan itu termasuk jenis yang menjadi dasar pembuatan diagram, biasanya basal, dan
bukan \omterm{def:b2:magmatic-differentiation:cumulate}{kumulat} atau lava yang telah berevolusi, yang \omterm{def:b2:elements-in-the-earth:trace}{unsur jejaknya} mencerminkan kristalisasi sama besarnya
dengan sumbernya. (2) Utamakan unsur yang tidak berpindah selama alterasi (Ti, Zr, Y, Nb, \omterm{def:b2:elements-in-the-earth:ree}{unsur tanah
jarang} berat) untuk batuan yang teralterasi atau termetamorfosis. (3) Plotlah serangkaian cuplikan, bukan satu:
satu titik di dekat batas tidak memutuskan apa pun. (4) Bacalah hasilnya sebagai kemungkinan:
medan-medan itu ditarik di sekeliling sebagian besar, bukan seluruh, cuplikan acuannya, dan beberapa medan
dipakai bersama. (5) Tegaskan dengan bukti independen: diagram lain, isotop,
konteks geologi.
\end{method}

Langkah (4) dapat dibuat kuantitatif (\cref{exo:b2:magmatic-series:8}). Jika sebuah basal jatuh di medan B, dan ketiga
tatanan sama-sama lazim, peluang bahwa basal itu MORB hanya sekitar 0.62.
Pengklasifikasi statistik yang dilatih dengan kompilasi semacam itu mewarisi batas yang sama: basal
membawa ingatan tentang sumber dan pelelehannya, dan dua tatanan dapat menghasilkan \omterm{def:b2:magma-emplacement:magma}{magma}
yang serupa.
% ledger: data:b4-basalt-settings

\begin{inthelab}[Memeriksa analisis batuan utuh]
Sebelum memplot batuan pada diagram mana pun, ahli petrologi memeriksa analisisnya. Oksida
utama, yang diukur misalnya dengan fluoresensi sinar-X, seharusnya berjumlah mendekati 100
bersama volatilnya: air dan karbon dioksida, yang diukur langsung atau bersama-sama sebagai
hilang pijar, yaitu massa yang hilang dari serbuk batuan ketika dipanaskan kuat. Hilang
pijar yang besar menandai batuan yang teralterasi, yang alkalinya mungkin telah berpindah; skema diagram
TAS menolak analisis dengan air lebih dari \qty{2}{\percent} atau
karbon dioksida lebih dari \qty{0.5}{\percent}. Analisis lalu dihitung ulang menjadi 100 tanpa
volatil, besi dinyatakan sebagai \ce{FeO} total atau dipisahkan menjadi fero dan feri,
dan bahan acuan yang dianalisis dalam pengukuran yang sama dibandingkan dengan nilai
tersertifikasinya (\cref{ch:b2:field-lab-methods-2}).
\end{inthelab}
% ledger: bgs:tas-use

\section{Latihan}

\begin{exercise}[$\star$]\label{exo:b2:magmatic-series:1}
Namailah batuan vulkanik berikut pada diagram TAS (\ce{SiO2}, \ce{Na2O}, \ce{K2O}
dalam \unit{\percent}, sudah tanpa volatil): (a) 49.60, 2.22, 0.51 (basal
acuan Hawaii BHVO-2); (b) 58.5, 3.9, 1.6; (c) 72.0, 4.0, 4.2; (d) 47.0, 4.5, 2.5.
Untuk (d), sebutkan dua batuan yang mungkin dan apa yang menentukan pilihan di antara keduanya.
\end{exercise}
% ledger: usgs:bhvo2-major, bgs:tas

\begin{exercise}[$\star$]\label{exo:b2:magmatic-series:2}
Sebuah lava memberikan \ce{SiO2} 51.60, \ce{TiO2} 1.20, \ce{Al2O3} 16.10, \ce{FeO_T} 9.80,
\ce{MgO} 6.60, \ce{CaO} 8.20, \ce{Na2O} 3.50, \ce{K2O} 1.40, \ce{H2O} 1.40, dan
\ce{CO2} \qty{0.20}{\percent}. (a) Apakah analisis itu dapat diterima untuk diagram TAS?
(b) Hitung ulang analisis itu tanpa volatil. (c) Namailah batuan itu sebelum dan sesudah
perhitungan ulang. (d) Apakah batuan itu bertipe K-rendah, K-sedang, atau K-tinggi?
\end{exercise}
% ledger: bgs:tas-use, bgs:k2o-sio2

\begin{exercise}[$\star$]\label{exo:b2:magmatic-series:3}
Namailah pada \omterm{def:b2:magmatic-series:qapf}{diagram QAPF}: (a) batuan berbutir kasar dengan \qty{25}{\percent} kuarsa,
\qty{30}{\percent} feldspar alkali, \qty{35}{\percent} plagioklas, \qty{8}{\percent}
biotit, dan \qty{2}{\percent} hornblende; (b) batuan dengan \qty{5}{\percent} kuarsa,
\qty{5}{\percent} feldspar alkali, \qty{55}{\percent} plagioklas, dan
\qty{35}{\percent} klinopiroksen.
\end{exercise}
% ledger: bgs:qapf

\begin{exercise}[$\star$]\label{exo:b2:magmatic-series:4}
Golongkan sebagai K-rendah, K-sedang, atau K-tinggi, dengan garis-garis pada gambar \ce{K2O}--\ce{SiO2}:
(a) 50.0, 0.40; (b) 56.0, 1.50; (c) 62.0, 2.90; (d) 66.0, 1.00 (\ce{SiO2},
\ce{K2O}, \unit{\percent}).
\end{exercise}
% ledger: bgs:k2o-sio2

\begin{exercise}[$\star\star$]\label{exo:b2:magmatic-series:5}
(a) Hitunglah $A$, $F$, dan $M$ untuk basal dengan \ce{FeO_T} 14.5, \ce{MgO} 4.5, dan
\ce{Na2O + K2O} \qty{3.2}{\percent}, dan untuk andesit dengan 6.1, 3.8, dan 5.05.
(b) Garis pemisah melalui $(F, A, M) = (54.5, 23.0, 22.5)$, $(52.5, 30.0,
17.5)$, dan $(49.6, 35.0, 15.4)$, dan titik pertama adalah titiknya yang paling kaya besi;
sisi toleiitik adalah sisi F. Di sisi mana setiap batuan jatuh? (c) Mineral
mana, yang mengkristal lebih awal, membuat andesit itu tetap berada di tempatnya?
\end{exercise}
% ledger: geoplotr:afm, vp21:series

\begin{exercise}[$\star\star$]\label{exo:b2:magmatic-series:6}
MORB median dalam kompilasi itu mengandung \qty{4}{ppm} Nb dan \qty{4.12}{ppm} La,
\omterm{def:b2:elements-in-the-earth:morb-oib}{basal pulau samudra} median 38 dan 30.8, dan basal busur median 3 dan 6.1. Dengan nilai
\omterm{def:b2:elements-in-the-earth:bse}{mantel primitif} Nb 0.658 dan La \qty{0.648}{ppm}, hitunglah nilai ternormalisasi
dan $\mathrm{Nb_N}/\mathrm{La_N}$ untuk masing-masing. Mana yang memiliki anomali niobium, dan mengapa?
\end{exercise}
% ledger: data:b4-basalt-settings, ms95:pm-trace

\begin{exercise}[$\star\star$]\label{exo:b2:magmatic-series:7}
Ti adalah $0.5993$ dari massa \ce{TiO2}. Hitunglah koordinat plot Ti--Zr--Y
(Ti/100, Zr, $3\mathrm Y$, dihitung ulang menjadi 100) untuk tiga basal: (a) \ce{TiO2}
\qty{1.6}{\percent}, Zr 110, Y \qty{32}{ppm}; (b) 2.8, 220, 29; (c) 0.8, 50, 20. Tentukan letaknya
pada gambar dan sebutkan kesimpulan apa yang dimungkinkan oleh setiap letak itu.
\end{exercise}
% ledger: geoplotr:tizry

\begin{exercise}[$\star\star$]\label{exo:b2:magmatic-series:8}
Dalam uji Ti--Zr--Y pada bab ini, medan B memuat 142 dari 164 MORB, 4 dari 173
\omterm{def:b2:elements-in-the-earth:morb-oib}{basal pulau samudra}, dan 101 dari 201 basal busur. (a) Berapa fraksi setiap kelompok yang jatuh
di B? (b) Jika ketiga tatanan sama-sama lazim, berapa peluang bahwa
basal yang jatuh di B adalah MORB? (c) Apa yang dapat dikatakan berdasarkan medan B saja?
\end{exercise}
% ledger: data:b4-basalt-settings

\begin{exercise}[$\star\star$]\label{exo:b2:magmatic-series:9}
Sebuah riolit mengandung \ce{Na2O} 5.5, \ce{K2O} 4.5, dan \ce{Al2O3} \qty{10.5}{\percent} ($M$ =
61.98, 94.20, dan \qty{101.96}{g/mol}). Hitunglah rasio molekuler $(\ce{Na2O +
K2O})/\ce{Al2O3}$-nya. (a) Apakah riolit itu peralkali? (b) Feldspar alkali mengandung satu
oksida alkali per \ce{Al2O3}: berapa fraksi alkali yang tidak dapat masuk ke dalamnya?
\end{exercise}
% ledger: bgs:tas-use

\begin{exercise}[$\star\star\star$]\label{exo:b2:magmatic-series:10}
Dalam model proposisi log-rasio, ambillah $\lambda_A = 0.4$ dan $\lambda_M =
0.8$. (a) Hitunglah kemiringan tren untuk $\lambda_F = 0.5$ dan untuk $\lambda_F =
0.7$. (b) Batuan Islandia memiliki kemiringan $-1.10$ dan batuan Cascade
$-0.72$: $\lambda_F$ berapa yang tersirat oleh masing-masing? (c) Seri mana yang kehilangan besinya lebih cepat,
dan mineral mana yang menjelaskannya? (d) Apa yang terjadi pada $F/M$ di sepanjang setiap tren?
\end{exercise}
% ledger: vp21:model

\begin{exercise}[$\star\star\star$]\label{exo:b2:magmatic-series:11}
Norma yang disederhanakan. Bekerjalah dalam mol oksida. Masukkan seluruh \ce{K2O} ke dalam ortoklas
(\ce{K2O.Al2O3.6SiO2}), seluruh \ce{Na2O} ke dalam albit (\ce{Na2O.Al2O3.6SiO2}),
\ce{Al2O3} yang tersisa bersama \ce{CaO} dalam jumlah yang sama ke dalam anortit (\ce{CaO.Al2O3.2SiO2}),
\ce{CaO} yang tersisa bersama $(\ce{Mg,Fe})\ce{O}$ dalam jumlah yang sama ke dalam diopsid
(\ce{CaO.(Mg,Fe)O.2SiO2}), dan sisa $(\ce{Mg,Fe})\ce{O}$ ke dalam hipersten
(\ce{(Mg,Fe)O.SiO2}). Jika silika kurang, ubahlah hipersten menjadi olivin
(\ce{2(Mg,Fe)O.SiO2}); jika masih kurang, ubahlah albit menjadi nefelin
(\ce{Na2O.Al2O3.2SiO2}). Terapkan pada (a) BHVO-2: \ce{SiO2} 49.60, \ce{Al2O3} 13.44,
\ce{FeO_T} 11.15, \ce{MgO} 7.257, \ce{CaO} 11.40, \ce{Na2O} 2.219, \ce{K2O}
\qty{0.513}{\percent}; (b) basal dengan 45.5, 15.0, 11.5, 8.5, 10.0, 3.8, dan 1.5 (massa
molar \ce{SiO2} 60.08, \ce{Al2O3} 101.96, \ce{FeO} 71.84, \ce{MgO} 40.30, \ce{CaO}
56.08, \ce{Na2O} 61.98, \ce{K2O} \qty{94.20}{g/mol}). Apakah masing-masing jenuh silika? Ke
seri mana setiap basal itu termasuk?
\end{exercise}
% ledger: usgs:bhvo2-major, bgs:series-terms

\begin{exercise}[$\star\star\star$]\label{exo:b2:magmatic-series:12}
Dalam kompilasi itu, $\mathrm{Nb}/\mathrm{La} < 0.7$ berlaku untuk 116 dari 141 basal busur,
21 dari 101 MORB, dan 3 dari 137 \omterm{def:b2:elements-in-the-earth:morb-oib}{basal pulau samudra} yang dianalisis untuk kedua unsur. (a) Hitunglah
ketiga fraksi itu. (b) Dengan proporsi ketiga tatanan yang sama, berapa
peluang bahwa basal dengan $\mathrm{Nb}/\mathrm{La} < 0.7$ adalah basal busur? (c)
Ulangi untuk wilayah tempat \qty{60}{\percent} basalnya MORB, \qty{30}{\percent}
\omterm{def:b2:elements-in-the-earth:morb-oib}{basal pulau samudra}, dan \qty{10}{\percent} basal busur. (d) Pelajaran apa yang dikandung (c)
bagi pemakaian \omterm{def:b2:magmatic-series:discrimination}{diagram diskriminasi}?
\end{exercise}
% ledger: data:b4-basalt-settings

\section{Soal: Satu gunung api, tiga lava}

\begin{problem}[{Soal akhir pekan --- dari tiga analisis sampai tatanan sebuah gunung api dan
kristal yang ditinggalkannya}]
\label{pb:b2:magmatic-series:1}
Tiga lava dari sebuah gunung api strato (rekaan), dari aliran lerengnya yang tertua sampai kubahnya
yang termuda, memberikan (\unit{\percent}):

\smallskip\noindent{\small
\begin{tabular}{@{}lrrrrrrrrrr@{}}
 & \ce{SiO2} & \ce{TiO2} & \ce{Al2O3} & \ce{FeO_T} & \ce{MgO} & \ce{CaO} & \ce{Na2O} & \ce{K2O} & \ce{H2O} & \ce{CO2}\\
lava 1 & 50.4 & 0.95 & 17.6 & 8.6 & 7.9 & 10.2 & 2.9 & 0.62 & 0.70 & 0.10\\
lava 2 & 58.6 & 0.80 & 17.4 & 6.1 & 3.8 & 6.9 & 3.9 & 1.15 & 1.10 & 0.05\\
lava 3 & 65.3 & 0.55 & 16.2 & 4.0 & 1.6 & 4.3 & 4.4 & 1.85 & 1.20 & 0.05\\
\end{tabular}}

\smallskip\noindent dan dalam ppm: Rb 10, 20, 32; Th 1.1, 2.1, 3.4; Nb 2.0, 3.6, 5.6;
La 7, 13, 20; Sr 520, 480, 380; Zr 70, 125, 190; Y 18, 22, 24. \omterm{def:b2:elements-in-the-earth:bse}{Mantel primitif}: Nb
0.658, La 0.648, Th \qty{0.0795}{ppm}.

\noindent\textbf{Bagian I --- Menamai lava.}
\begin{enumerate}
\item Jumlahkan setiap analisis. Apakah ketiganya dapat diterima untuk diagram TAS?
\item Hitung ulang lava 1 tanpa volatil: berikan faktornya, \ce{SiO2}, dan
\ce{Na2O + K2O}.
\item Lakukan hal yang sama untuk lava 2 dan 3.
\item Namailah ketiga lava itu pada diagram TAS.
\item Apakah namanya akan berubah tanpa perhitungan ulang?
\item Sebuah batuan berbutir kasar yang berkomposisi seperti lava 2 mengkristal di akar
gunung api itu. Skema mana yang menamainya, dan di medan mana kamu mengharapkannya?
\end{enumerate}

\noindent\textbf{Bagian II --- Seri yang mana?}
\begin{enumerate}[resume]
\item Golongkan setiap lava sebagai K-rendah, K-sedang, atau K-tinggi.
\item Hitunglah $A$, $F$, $M$ untuk masing-masing (tanpa volatil).
\item Garis pemisah melalui $(F, A, M) = (42.8, 17.1, 40.1)$, $(46.9, 18.1,
35.0)$, $(52.5, 30.0, 17.5)$, $(49.6, 35.0, 15.4)$, $(40.8, 50.0, 9.2)$, dan $(37.9,
55.0, 7.1)$, dengan medan toleiitik di sisi F. Dengan menginterpolasi $F$ secara linear terhadap
$A$ di antara titik-titik yang bertetangga, di sisi mana setiap lava jatuh?
\item Hitunglah \ce{FeO_T}/\ce{MgO} untuk setiap lava. Apakah ada pengayaan besi?
\item Ke seri mana gunung api itu termasuk? Mineral mana dan kondisi apa yang
menjelaskannya?
\end{enumerate}

\noindent\textbf{Bagian III --- Tatanan yang mana?}
\begin{enumerate}[resume]
\item Normalkan Nb, La, dan Th lava 1 terhadap \omterm{def:b2:elements-in-the-earth:bse}{mantel primitif}.
\item Hitunglah $\mathrm{Nb_N}/\mathrm{La_N}$ untuk lava 1 dan bandingkan dengan nilai-nilai median
pada \cref{exo:b2:magmatic-series:6}.
\item Hitunglah $\mathrm{Th_N}/\mathrm{Nb_N}$. Apa yang membawa torium itu?
\item Hitunglah koordinat Ti--Zr--Y lava 1 (\ce{TiO2} dari analisis tanpa
volatil). Apa yang dinyatakan oleh medannya? Mengapa lava 2 dan 3 tidak diplot?
\item Simpulkan: di mana gunung api ini tumbuh?
\end{enumerate}

\noindent\textbf{Bagian IV --- Berapa banyak yang mengkristal?}
\begin{enumerate}[resume]
\item Kalium oksida tidak masuk ke mineral awal mana pun. Dari \ce{K2O} tanpa volatil
lava 1 dan 2, hitunglah fraksi $F$ cairan yang tersisa ketika lava 2 terbentuk, jika lava itu
berasal dari lava 1 melalui \omterm{def:b2:magmatic-differentiation:crystallisation}{kristalisasi fraksional}.
\item Lakukan hal yang sama dengan rubidium dan dengan torium.
\item Zirkonium memiliki koefisien total $D \approx 0.1$. Gunakan hukum Rayleigh untuk menaksir $F$ darinya.
\item Stronsium turun dari lava 2 ke lava 3. Hitunglah $F$ dari kalium oksida antara lava 2 dan
3, lalu gunakan untuk menaksir $D$ total Sr pada selang itu. Mineral mana yang
bertanggung jawab?
\item Dari lava 1 sampai lava 3, berapa fraksi cairan yang telah mengkristal?
\item Seandainya \omterm{def:b2:magma-emplacement:magma}{magma} itu mengasimilasi kerak granitik yang kaya akan kalium oksida, apakah taksiran
dari kalium oksida akan terlalu tinggi atau terlalu rendah?
\item Rata-ratakan keempat taksiran Bagian IV untuk lava 1 sampai 2 dan nyatakan massa
kristal yang tertinggal per \qty{100}{kg} basal untuk menghasilkan andesit itu, yaitu angka kunci
soal ini.
\end{enumerate}
\end{problem}
% ledger: bgs:tas-use, bgs:tas, bgs:k2o-sio2, geoplotr:afm, geoplotr:tizry, ms95:pm-trace
